% !TeX program = xelatex
% !TeX program = xelatex
% ===========================================================================
%  tech_common.tex -- 两份技术文档共用的导言区
% ===========================================================================
\documentclass[11pt,a4paper]{ctexart}

\usepackage[top=2.4cm,bottom=2.4cm,left=2.4cm,right=2.4cm]{geometry}
\usepackage{amsmath,amssymb,amsthm}
\usepackage{graphicx}
\usepackage{booktabs}
\usepackage{longtable}
\usepackage{array}
\usepackage{enumitem}
\usepackage{xcolor}
\usepackage{listings}
\usepackage{tcolorbox}
\usepackage{caption}
\usepackage{fancyhdr}
\usepackage{titlesec}
\usepackage{hyperref}

\hypersetup{colorlinks=true,linkcolor=blue!55!black,citecolor=blue!55!black,
            urlcolor=blue!55!black,bookmarksnumbered=true}

% ---------------------------------------------------------------- 字体/版式
\setlength{\emergencystretch}{2.5em}
\setlength{\parskip}{0.35em}
\setlength{\parindent}{2em}
\linespread{1.12}

% ---------------------------------------------------------------- 代码样式
\definecolor{codebg}{RGB}{247,247,250}
\definecolor{codekw}{RGB}{0,90,160}
\definecolor{codecm}{RGB}{0,128,96}
\definecolor{codest}{RGB}{170,60,60}
\lstset{
  basicstyle=\ttfamily\footnotesize,
  backgroundcolor=\color{codebg},
  frame=single, framerule=0.4pt, rulecolor=\color{gray!45},
  breaklines=true, breakatwhitespace=false,
  numbers=left, numberstyle=\tiny\color{gray!80}, numbersep=6pt,
  keywordstyle=\color{codekw}\bfseries, commentstyle=\color{codecm}\itshape,
  stringstyle=\color{codest}, showstringspaces=false, columns=flexible,
  tabsize=2, extendedchars=true, captionpos=b,
  xleftmargin=1.6em, framexleftmargin=1.4em
}
\renewcommand{\lstlistingname}{代码}

% ---------------------------------------------------------------- 提示框
\newtcolorbox{keybox}[1][]{colback=blue!4,colframe=blue!45!black,
  boxrule=0.6pt,arc=2pt,left=5pt,right=5pt,top=4pt,bottom=4pt,
  fonttitle=\bfseries,#1}
\newtcolorbox{warnbox}[1][]{colback=orange!6,colframe=orange!55!black,
  boxrule=0.6pt,arc=2pt,left=5pt,right=5pt,top=4pt,bottom=4pt,
  fonttitle=\bfseries,#1}

% ---------------------------------------------------------------- 定理环境
\newtheorem{theorem}{定理}
\newtheorem{proposition}{命题}
\newtheorem{lemma}{引理}
\renewcommand{\proofname}{\heiti 证明}

% ---------------------------------------------------------------- 页眉页脚
\pagestyle{fancy}
\fancyhf{}
\fancyhead[L]{\small\color{gray!70}\leftmark}
\fancyhead[R]{\small\color{gray!70}\thepage}
\renewcommand{\headrulewidth}{0.3pt}
\renewcommand{\headrule}{\hbox to\headwidth{\color{gray!40}\leaders\hrule height \headrulewidth\hfill}}

\titleformat{\section}{\normalfont\heiti\Large}{\thesection}{0.6em}{}
\titleformat{\subsection}{\normalfont\heiti\large}{\thesubsection}{0.6em}{}
\titleformat{\subsubsection}{\normalfont\heiti\normalsize}{\thesubsubsection}{0.6em}{}

\captionsetup{font=small,labelfont=bf,labelsep=quad}


\title{\heiti 可选优化尝试与负优化报告\\[4pt]
       \large 针对 §10.2 四条改进方向的实测结论\\[2pt]
       \normalsize —— 11 项尝试中 6 项为负优化、2 项无效、3 项有效}
\author{2026 高教社杯全国大学生数学建模竞赛 B 题}
\date{\today}

\begin{document}
\maketitle
\thispagestyle{fancy}

\begin{abstract}
\noindent
论文 §10.2 列出了四个改进方向。本文档把它们逐一做成\textbf{可开关的实现}，
在\textbf{同一个算例池}上对照实测，并如实记录每一项的结果——\textbf{包括做得更差的那些}。

\textbf{结论概览}：共尝试 12 项改动。\textbf{4 项有效}（朝估计点归航、垂向截断测量、自适应沿射线站位、自适应补测间隔），\textbf{6 项为负优化}（边境外补测、提高清除预算的三组参数、更大的末段半径、更密的格点），\textbf{2 项无效}（顺路清除耦合、progress\_ratio）。

其中最重要的两项有效改动直接解决了 §10.2 第 1 条：\\
\textbf{垂向截断测量}——末段三次观测点近乎共线导致 $\sigma$ 停在 77 m，
从垂直于视线方向补一次观测后 $\sigma$ 降到 \textbf{1 m}；\\
\textbf{自适应沿射线站位}——贴边朝外的定向源其受光区是薄月牙，
原来的 $\pm55^{\circ}$ 偏移站位直接落到背光侧，改为随失败次数收缩偏移角、最终回到射线上。
连同\textbf{朝估计点归航}的修正，问题四在\textbf{独立算例池}上的完成率由 $0.9981$ 升到 \textbf{$1.0000$}（最差由 $0.9231$ 升到 $1.0000$），同时平均时间下降 \textbf{$36.7\%$}。

至于 §10.2 的第 3、4 条，本文用实验判定其\textbf{必要性}：
把示向度误差换成\textbf{空间相关}的平滑随机场（相关比例 $0/50/80/100\%$）后结果无系统性变化，
说明本文的\textbf{最坏情形几何}估计不依赖误差独立性；
真实耗时方面，一次请求实测 $6.5$ ms，距离 1200 s 预算有 \textbf{232 倍}余量。
因此这两条\textbf{无需建模}，但第 4 条在官方 GUI 模拟器上仍需实测确认。
\end{abstract}

\tableofcontents
\newpage

% ===========================================================================
\section{方法与判据}

\subsection{怎么算"变好了"}

每一项改动都做成\textbf{可开关的参数}（默认关闭），在\textbf{同一个算例池}上对照。
比较采用词典序，并把两侧都记录下来：

\begin{enumerate}[leftmargin=1.8em]
  \item 完成率（越高越好，这是第一优先级）；
  \item 最差算例完成率；
  \item 平均定位清除时间（越低越好）；
  \item 平均行程（解释时间变化的物理量）。
\end{enumerate}

\textbf{判定规则}：完成率下降 $>0.002$ 判为负优化；完成率持平而时间上升 $>2\%$ 判为负优化；
完成率持平而时间下降 $>2\%$ 判为有效；其余判为无效（噪声级）。

\begin{warnbox}[title=为什么必须把负结果写下来]
上一轮已经吃过一次教训：一次"下降 44\%"的调参结论在独立算例池上只剩 1.7\%，
根因是算例池被复用而 \texttt{Arena} 会就地清除 \texttt{Source} 对象。
\textbf{只要对照实验本身是脏的，正结果和负结果一样不可信}。
本轮所有对照都在同一批固定算例上、每次重新生成算例、并且同一配置连算两次必须同分。
\end{warnbox}

% ===========================================================================
\section{§10.2 第 1 条：贴边朝外源的完成率损失}

\subsection{先诊断，不要先动手}

一开始的假设是"这些源在域内不可认证，所以要到域外采样"（即上一版实现的 rim patrol）。
诊断脚本 \texttt{diagnose\_loss.py} 给出的画面却完全不同：

\begin{lstlisting}[caption={失败算例的真实状态（seed 30017，频道 11）}]
true: pos=(92,1409) r=1412 m  kind=directional  dir=111 deg  R_eff=1007 m
knowledge: status=located   n_obs=5   est=(111, 1401, 77)
last task counters: attempts=6 total=6 vantage=3
measured positions inside the reception radius: 13 / 34
... of which on the LIT side: 5
\end{lstlisting}

三点关键事实：\textbf{(i)} 源\textbf{已经被定位}，估计与真值只差 $20.1$ m；
\textbf{(ii)} 但 $\sigma$ 报告为 77 m，末段六边形图案的半径取 $0.55\sigma\approx42$ m，
图案点全在 42 m 与 85 m 两个圈上，\textbf{没有一个落在 20 m 清除半径内}；
\textbf{(iii)} 逐迭代轨迹显示，图案扫完后机器人\textbf{停在最后一个图案点上}
$(184,1359)$——距离真值 104 m，而且\textbf{在定向源的背光侧}，
于是随后的"原地再听一次"必然返回 no\_signal，整个循环空转到预算耗尽。

\subsection{有效改动 1：图案把中心留到最后}

\begin{keybox}[title=改动]
把末段六边形清除图案的尝试顺序改为\textbf{先扫外圈、最后打中心}。
这样一次失败的图案扫描结束时，机器人正好站在\textbf{估计点}上——也就是最值得再听一次的位置；
原来的顺序会让它停在最后一个外圈点上，可能位于背光侧。\textbf{代价为零}（同样的动作数，只是顺序不同）。
\end{keybox}

\noindent\textbf{效果}：静止位置由 $(184,1359)$ 变为 $(111,1401)$，即估计点。

\subsection{有效改动 2：垂向截断测量（本轮最大的单项收益）}

即便站到了估计点上，问题仍未解决：估计点落在\textbf{背光侧}
（与真值的内积 $-14.4<0$），在那里听不到；而估计点距离真值 $20.5$ m，
\textbf{恰好比 20 m 的清除半径多 0.5 m}。

根因在更上游：\textbf{三次观测点几乎共线}。归航是沿直线走的，
三个观测点 $(-288,1538)$、$(-42,1455)$、$(38,1428)$ 近似在一条直线上，
它们的角楔交会得到一个\textbf{细长条}，最小包围圆半径因此停在 77 m——
\textbf{共线的方位测量无法约束距离}，这是纯方位定位的经典性质。

\begin{keybox}[title=改动]
当末段图案失败且 $\sigma$ 仍大于阈值时，从\textbf{垂直于当前视线}的方向
偏移约 80 m 再测一次，然后重估。共线的三次切割约束不了距离，一次垂向切割就能约束。
代价：一次检测加约 80 m 行程（约 20 s），而此前机器人为一个"其实已经定位到"的频道多搜了 1000 s 以上。
\end{keybox}

实测轨迹（同一算例）：

\begin{lstlisting}[caption={垂向截断前后}]
改动前: est=(111,1401,77) -> 图案全miss -> 原地听不到 -> 预算耗尽 -> 失败 (1578 s)
改动后: est=(111,1401,77) -> 垂向走 54 m 到 (108,1432) -> 听到 direction
        -> new est=(92,1409,1)      <-- sigma 从 77 m 降到 1 m
        -> 清除成功, 距真值 11.6 m
该案例: ratio 0.933 -> 1.000, 平均时间 1577.9 s -> 739.1 s (-53%)
\end{lstlisting}

\subsection{有效改动 3：自适应沿射线站位}

独立算例池上仍有两个失败算例，诊断显示两者都是\textbf{贴在边界上、且定向方向朝外}的源：

\begin{center}
\begin{tabular}{lcccc}
\toprule
算例 & 真实位置 & 半径 & 定向方向 & 失败时状态 \\ \midrule
seed 62016 & $(-753,-1618)$ & 1784 m & $224^{\circ}$ & 已定位，估计误差 20.8 m \\
seed 62029 & $(288,1668)$ & 1692 m & $93^{\circ}$ & \textbf{只听到 1 次}，站位预算耗尽 \\
\bottomrule
\end{tabular}
\end{center}

对于这类源，\textbf{域内的受光区只是一个薄月牙}。问题二的站位规则把第二观测点放在
离首观测点约 550 m、偏离示向度 $\pm55^{\circ}$ 的位置——这个点\textbf{直接落到背光侧}，
测量返回 no\_signal，三次尝试全部浪费（seed 62029 正是如此）。

\begin{keybox}[title=改动]
站位偏移角随失败次数\textbf{自适应收缩}：第 1 次 $\le45^{\circ}$，
第 2 次 $\le20^{\circ}$，第 3 次\textbf{回到射线上}（偏移 $0^{\circ}$，距离取 $0.55$ 倍可行射线长度）。
依据是\textbf{射线逼近引理}：从曾听到信号的点出发、指向源的射线上任意一点都在辐射扇区内，
因此射线上的站位\textbf{保证能听到}；代价是几何交会角变差。
贴边场景下"能不能听到"是硬约束，交会角只是次要因素。
\end{keybox}

\subsection{有效改动 4：沿估计方向归航（本节最大的单项收益）}

在把前两项改动落地后，正式测试里仍有一个算例失败。追踪它时发现了\textbf{一个更基础的缺陷}：

\begin{lstlisting}[caption={案例 seed 910001、频道 18 的归航轨迹（修复前）}]
target channel 18 at (-1291.0, -445.1)  directional  dir=311 deg  R_eff=1448 m
>>> clear_channel(18) call #1  pos=(-1639.5,-178.0)  dist to TRUE source = 439.1 m
    estimate (-1287.4, -443.0) sigma=111.1  est error = 4.2 m      <-- 已定位到 4.2 m!
   [it1] walk 420 m -> (-818,  -43)     <-- 朝反方向走
   [it2] walk 420 m -> (-1140,-313)     <-- 走回来
   [it3] walk 189 m -> (-963, -164)     <-- 又走开
   ...
RESULT ratio=0.909   mean=1732.9 s   travel=65854 m
\end{lstlisting}

\textbf{机器人已经把一个源定位到 4.2 m（清除半径是 20 m），却在离它 400 多米的地方来回振荡。}

根因在归航步的方向计算上。原实现取\textbf{最后一条观测的方位角}，再从这个方位角
\textbf{从当前位置}前进——但方位角只在\textbf{测量它的那个点}上有意义。
本例中最后一条方位是在很远的位置测的，两者方向的夹角达 $58^{\circ}$，
于是每一步都把机器人送得更远；步长又按"过冲"规则逐步减半，
形成"走远—回退—再走远"的极限环，直到尝试预算耗尽。

\begin{keybox}[title=改动]
归航步改为\textbf{朝估计点前进}：$b=\operatorname{atan2}(\hat q_y-p_y,\ \hat q_x-p_x)$。
当机器人确实站在最后观测点附近时，两个方向本来就重合，因此没有任何损失；
当它不在时，朝估计点前进是唯一有意义的方向。
\end{keybox}

修复后同一算例：\textbf{ratio 0.909 $\to$ 1.000，平均时间 1732.9 s $\to$ 1133.1 s（$-35\%$）}，
归航稳定收敛（估计误差 4.0 m、$\sigma$ 由 18.9 降到 12，最终在 9.9 m 处清除）。

\begin{warnbox}[title=这条缺陷为什么值得单独记下来]
它不是被"更强的算法"修好的，而是被\textbf{把每一步的位置和目标打印出来}发现的。
三种改动（垂向截断、沿射线站位、朝估计点前进）里，最后这一条实现最简单、收益最大，
却最不容易靠读代码发现——因为"用最后一条方位前进"在\textbf{正常情况下}是对的
（机器人总是先走到观测点再前进），只有在异常路径上才暴露。
\textbf{任何迭代式控制都必须把"当前状态与判据"打出来看。}
\end{warnbox}

\subsection{第 1 条的最终实测}

在\textbf{从未参与过任何调参}的独立算例池（种子 62000--62039，40 例）上：

\begin{center}
\begin{tabular}{lcccc}
\toprule
& 完成率 & 最差完成率 & 平均时间/s & 平均行程/m \\ \midrule
优化前 & 0.9940 & 0.9091 & 1382.9 & 68925 \\
优化后 & \textbf{1.0000} & \textbf{1.0000} & \textbf{901.0} & \textbf{45125} \\
变化 & $+0.0060$ & $+0.0909$ & $\mathbf{-34.8\%}$ & $-34.5\%$ \\
\bottomrule
\end{tabular}
\end{center}

\textbf{§10.2 第 1 条所列的"约 4\% 完成率损失"已消除}
（40 例全部 100\%），而且时间同时下降约三分之一。
值得注意的是：真正解决问题的\textbf{不是}域外采样，而是"把定位做准"与"把站位放对"。

% ===========================================================================
\section{§10.2 第 2 条：利用清除顺序与探测的耦合}

\subsection{两个具体化的耦合}

\begin{itemize}[leftmargin=1.8em]
  \item \textbf{C1 顺路清除}：巡访搜索站位时，顺带清除半径 260 m 内已经定位但未清除的源
        ——行程已经付过，边际成本几乎为零。
  \item \textbf{C2 自适应补测间隔}：机会式补测的间隔固定为 650 m；
        当未解算频道仍多于 5 个时信息更值钱，间隔收缩到 320 m。
\end{itemize}

\subsection{实测（24 例，同一算例池）}

\begin{center}
\begin{tabular}{lcccc}
\toprule
方案 & 完成率 & 平均时间/s & 平均行程/m & 判定 \\ \midrule
基线（无耦合） & 1.0000 & 881.3 & 44356 & --- \\
C1 顺路清除 & 1.0000 & 882.2 & 44407 & \textbf{无效}（差 0.1\%，噪声级） \\
C2 自适应补测 & 1.0000 & \textbf{852.1} & \textbf{41369} & 有效（$-3.3\%$ 时间，$-6.7\%$ 行程） \\
C1 $+$ C2 & 1.0000 & 851.6 & 41277 & 与 C2 相同（C1 无增量） \\
\bottomrule
\end{tabular}
\end{center}

\begin{warnbox}[title=一个必须如实报告的负结果]
\textbf{C1 无效}。原因事后看是清楚的：任务序列已经为清除任务加了 260 s 的\textbf{优先权偏置}，
最近邻排序时清除任务本来就会被提前访问，所以"顺路清除"这一显式耦合没有增量。
这也说明论文原先关于"未利用清除顺序与探测的耦合"的自评\textbf{偏保守}：
耦合的一部分已经被优先权代价隐含地利用了。C1 已默认关闭并保留为记录。
\end{warnbox}

% ===========================================================================
\section{§10.2 第 3 条：误差空间相关性}

\subsection{第一次实验的设计是错的}

第一轮实验给出"$a=50\%$ 时完成率下降、$80\%/100\%$ 反而正常"的怪现象。
复核后发现\textbf{三种误差的分布不可比}：独立误差是 $U(-1,1)$，标准差 $2/\sqrt{12}=0.577$；
而三个正弦叠加的平滑场标准差只有 0.419。也就是说\textbf{"100\% 相关"其实是一个更容易的问题}，
$50\%$ 那一档正处在"误差幅度被压低、但还没有获得平滑场的小幅度"的中间位置
——它看起来最差，只是因为它既不是原来的难度、也不是更低的难度。

把平滑场\textbf{按标准差归一化}（缩放因子 1.378）后重做，才是公平对照：

\begin{center}
\small
\begin{tabular}{lccccc}
\toprule
$a$ & 标准差 & 截断比例 & 50 m 误差差 & 漏源 / 总数 & 完成率 \\ \midrule
0\%  & 0.577 & 0.0\% & 0.834 & \textbf{2--3 / 914}（两轮独立复测） & 0.9969$\sim$0.9979 \\
50\% & 0.407 & 0.7\% & 0.419 & \textbf{2--4 / 914} & 0.9951$\sim$0.9979 \\
100\%& 0.551 & 8.8\% & 0.088 & \textbf{4 / 914} & 0.9952 \\
\bottomrule
\end{tabular}
\end{center}

\subsection{结论：相关性的影响真实存在，但在本问题里是二阶效应}

\begin{enumerate}[leftmargin=1.8em]
  \item \textbf{原先观察到的"50\% 更差"是噪声}：$a=0$（即当前实现）在同样的池子上也会漏
        2$\sim$3 个源。914 个源里漏 2 个与漏 4 个的差别在统计上不可分辨
        （要可靠区分 2 倍差异大约需要 10 倍样本量）。
  \item \textbf{相关性确实带来一点损失}：公平设计下 $a=100\%$ 的漏源数约为 $a=0$ 的两倍
        （$0.22\%\to0.44\%$）。机制清楚：\textbf{共享偏差不会在求交时相互抵消}——
        邻域内所有方位被同向偏转，交会区域被整体平移；而独立误差在多个约束求交时会部分抵消。
        这也说明现实中"在同一片区域多测几次"的收益比独立误差假设下要小。
  \item \textbf{但本文的估计器对此免疫}：定位区域取的是"每条读数都落在
        $\hat\theta\pm1^{\circ}$ 内"这一\textbf{硬约束的交集}，结论对任意满足逐条界的误差实现都成立。
        相关性只影响"多余读数还能收紧多少"，而本文的策略并不依赖这一点——
        \textbf{重测的价值来自交会角的改善，而不是误差的平均}。
\end{enumerate}

\subsection{顺带暴露出来的真正残留失败模式}

公平对照把一个此前被掩盖的事实暴露出来：\textbf{即使误差完全独立，仍有
$0.2\sim0.35\%$ 的源被漏掉}（约 900 个源中漏 2$\sim$3 个）。把这些算例逐个打开，
它们的结构完全相同：

\begin{center}
\begin{tabular}{llcccc}
\toprule
算例 & 频道 & 真实位置 & 半径 & 类型 / 朝向 & 唯一一条方位的测量点 \\ \midrule
seed 100006 & ch 2 & $(-1760,-76)$ & 1761 m & 定向 $163^{\circ}$ & $(-1800,0)$（域边界上） \\
seed 102009 & ch13 & $(1717,21)$  & 1717 m & 定向 $356^{\circ}$ & $(1800,0)$（域边界上） \\
\bottomrule
\end{tabular}
\end{center}

共同点：\textbf{贴边（$r\approx1720\sim1760$ m）、定向、且唯一一条方位是在目标域边界上测到的}。
机器人当时离源只有 85 m、且位于受光侧，却因为下面这条推理链而永远停在这一步：

\begin{enumerate}[leftmargin=1.8em]
  \item 只有一条方位 $\Rightarrow$ 无法定位 $\Rightarrow$ 不生成清除任务；
  \item 站位规则把第二观测点放在沿射线 550 m 处——而该处已在\textbf{源之外 465 m 的背光侧}
        （射线逼近引理只保证"从听到点朝向源"的那一段受光，越过源之后不再保证）；
  \item 三次站位全部落在背光侧，返回 no\_signal，预算耗尽，该频道永久停留在"已听到"状态。
\end{enumerate}

针对这一点做了\textbf{两次}改进尝试，结果如下（均在同一批固定池子上对照）。

\subsubsection{尝试一：沿射线小步爬行（保留，中性）}

把沿射线的 550 m 单跳改为\textbf{200 m 的爬行}：丢失信号即判定"越过源"，
退回听到点并把步长乘 0.45。效果是\textbf{确实拿到了第二条方位}（观测数 $1\to2$），
但估计的 $\sigma$ 仍高达 670 m / 26180 m——因为\textbf{两条方位几乎平行}：

\begin{lstlisting}[caption={爬行后的实测（seed 100006）}]
obs0 at (-1800,   0) brg=298.5  dist_to_src=86 m  LIT
obs1 at (-1781, -36) brg=297.2  dist_to_src=45 m  LIT   <-- 只差 1.3 度
estimate (-1780.7, -35.6) sigma = 670 m   (真值 (-1760,-76), 位置其实只差 41 m)
\end{lstlisting}

爬行解决了"保持受光"，但没有解决"有基线"。在整体池子上它是\textbf{中性}的
（漏源 2$\sim$3/914 与改动前相同），因此\textbf{保留}（它至少让观测数从 1 变成 2，
对后续任何改进都是前提）。

\subsubsection{尝试二：短基线横向偏移（\textbf{负优化，已撤销}）}

观察到"已有一条近距方位、源就在 45 m 外"之后，尝试在第二条方位的基础上加一次
\textbf{35 m 的横向偏移}（角度放宽到 $70^{\circ}$）。直觉是：源这么近，35 m 的基线
在 45 m 距离上张角约 $40^{\circ}$，几何上足够，而且短到不会越出受光半平面。

实测（三个池子、914 个源、相同种子）：

\begin{center}
\begin{tabular}{lcccc}
\toprule
方案 & 漏源 / 总数 & 加权完成率 & 平均时间/s & 异常现象 \\ \midrule
仅爬行 & \textbf{2 / 914}（0.219\%） & 0.9979 & $\sim$890 & --- \\
爬行 $+$ 短基线 & \textbf{8 / 914}（0.875\%） & 0.9917 & $\sim$1050 & 单频道 5974 次 no\_signal \\
\bottomrule
\end{tabular}
\end{center}

\textbf{损失扩大 4 倍、时间增加 18\%，并且出现病态行为。}
原因是放宽后的预算条件让"已有两条方位"的频道可以无限申请站位，
它们挤占了本应用于普查覆盖的预算，却没有真正收紧 $\sigma$。
该改动已\textbf{撤销}，参数保留在代码里并标注为"实测负优化、未使用"。

\begin{warnbox}[title=这一节的三个可复用教训]
\textbf{(1) 对照实验本身要先检查是否公平}：两次实验的误差分布标准差不同（0.577 vs 0.419），
"更相关"其实是"更容易"。\textbf{改一个因素之前，先确认其他因素没被一起改掉。}\par
\textbf{(2) 单个算例的"修好了"不代表方向正确}：短基线改动在 seed 102009 上从失败变为成功，
在 seed 100006 上从失败变得更失败。\textbf{必须回到池子上看净效果。}\par
\textbf{(3) 预算型放宽几乎总是负优化}：本轮 N2--N6 与这里的尝试二都是同一模式
——"让机器人再试一次"只会让它在同一个错误方向上花更多时间。
\end{warnbox}

\section{§10.2 第 4 条：真实时间耦合——\textbf{在本地环境判定为不需要}}

测量一次完整的正式测试（问题四）的\textbf{真实耗时}：

\begin{center}
\begin{tabular}{lcc}
\toprule
量 & 数值 & 说明 \\ \midrule
请求次数 & 799 & 走完整 HTTP 协议 \\
HTTP 总耗时 & 5.18 s & 客户端侧累计 \\
每请求平均 & \textbf{6.48 ms} & \\
模拟器内部耗时 & 19.63 s & 含案例生成与日志落盘 \\
程序运行时间上限 & 1200 s & 题目规定 20 min \\
\textbf{余量倍数} & \textbf{232$\times$} & 需要每请求慢到 1.5 s 才会触限 \\
\bottomrule
\end{tabular}
\end{center}

\subsection{判定：不需要把真实时间纳入目标函数}

即使把仿真环境内部开销全部计入（19.63 s），余量仍有 61 倍。
用真实时间与虚拟时间做双目标优化，在当前环境下只会增加模型复杂度而不改变任何决策。

\begin{warnbox}[title=一个不能忽略的保留意见]
上述结论\textbf{只对本地模拟器成立}。官方模拟器是一个 Windows \textbf{GUI 程序}，
其响应时间由界面刷新与动画驱动，可能比本地回环服务慢得多。
本队无法获取官方模拟器（原因见自检报告 §2），因此\textbf{无法实测}。
可以给出的定量说法是：官方的每请求响应需要比本地慢 \textbf{230 倍以上}
才会让"程序运行时间"成为约束。若将来在官方环境上运行，这一项应重新评估——
这也是本文把该条保留在 §10.2 而不是删除的原因。
\end{warnbox}

% ===========================================================================
\section{还剩哪些优化方向？}

完成率已饱和（独立池 1.0000、演练池 60 例 1.0000、正式测试三局全清），
唯一还在动的指标是时间。先把时间花在哪里量清楚，再谈方向。

\subsection{时间预算实测（问题四，16 例均值）}

\begin{center}
\small
\begin{tabular}{p{3.4cm}p{1.8cm}p{2.2cm}p{2.0cm}p{2.4cm}}
\toprule
任务类 & 次数/例 & 行程/m & 占行程 & 占虚拟时间 \\ \midrule
普查站位（probe） & 25.5 & 29442 & 73.6\% & \textbf{67.6\%} \\
清除（clear） & 12.7 & 8895 & 22.2\% & 25.4\% \\
补测站位（localise） & 4.9 & 2476 & 6.2\% & 7.3\% \\
入场普查（census） & 1.0 & 0 & 0\% & 1.1\% \\
\bottomrule
\end{tabular}
\end{center}

按动作类型拆：\textbf{移动 75.2\%}、检测 20.1\%、频道切换 3.8\%、清除脉冲 0.9\%。
每例 429 次检测（约 2144 s）、6461 次切换、25.5 个搜索站位。
\textbf{结论：一切优化都必须打在"移动"上；检测与切换加起来只有 24\%。}

\subsection{方向一：搜索站位的路径（潜在收益最大，约 15$\sim$20\%）}

实际到访的 25.5 个站位，若按最近邻顺序从原点连起来只需 \textbf{18981 m}，
而实际行程是 40027 m。把"清除/补测"任务也算进去之后仍有明显余量：
搜索任务的行程 29442 m 对 18981 m 的巡回，\textbf{多出约 10.5 km，折合 2100 s
（约占总时间的 20\%）}。

原因是任务分两类（搜索与清除）被混在同一个滚动时域序列里反复重规划，
机器人会为清除任务离开普查路线再折返。可行的改法是把两者\textbf{分阶段}：
先按一次完整巡回走完全部普查站位，再集中清除；只有当"已知源位于巡回路线附近"
时才插队。这是一个结构性改动，需要重写任务调度而不是调参。

\subsection{方向二：每个站位的检测频道数（潜在收益约 10\%）}

每次到达站位都要对\textbf{全部未解算频道}做一遍检测（最多 20 个频道 $\times$ 6 s = 120 s）。
实际上认证只要求"候选位置周围有若干方向的读数"，
\textbf{远离该站位的候选点测了也没用}。按"本站位能推进哪些候选点的认证"来选频道，
预计可把 429 次检测压到 250 次左右，即约 1000 s/例（10\%）。

\subsection{方向三：残留的 0.2$\sim$0.35\% 漏源（完成率，非时间）}

见上一节：贴边定向源、唯一方位取自域边界。两次尝试（爬行、短基线）分别是中性与负优化。
真正要解决它需要\textbf{主动重捕获策略}：一旦某频道"已听到但未定位"，
就应该\textbf{回到那个听到点}，并沿"垂直于该方位"的方向做一系列\textbf{短步}探测
（步长由"还能听到多远"二分确定），直到取得一条方向差足够大的第二方位。
这与现有的"站位—清除"框架是并行的另一套机制，属于下一步工作。

\subsection{已经不剩下的方向}

\begin{itemize}[leftmargin=1.8em]
  \item \textbf{完成率}：在 914 个源的池子上已到 $0.997\sim0.998$，剩余部分就是方向三；
        在本文报告的全部演练与正式测试池上均为 1.0000。
  \item \textbf{检测与切换}：合计 24\%，且已经按频道号升序测量（切换次数已接近下界）。
  \item \textbf{清除动作}：清除脉冲只占 0.9\%，图案化清除已经把"定位"转成了"清除"。
  \item \textbf{参数}：8$\sim$9 个参数在 313$\sim$348 s（问题三）与 850$\sim$1050 s（问题四）
        的窄带内无显著差异，已到平台。
\end{itemize}

\section{负优化清单（完整记录）}

下列改动\textbf{都做了实现、都做了对照实验，结果都是更差或无效}。
把它们写下来，是为了让后来者不必重复。

\footnotesize
\begin{longtable}{p{0.6cm}p{3.4cm}p{3.1cm}p{3.0cm}p{2.4cm}}
\toprule
\\序号 & 尝试 & 结果 & 判定 & 依据 \\ \midrule
\endhead
N1 & 域外补测环 / rim patrol（原 §10.2 第 1 条的方案） &
完成率完全相同（0.9944，最差 0.9000），时间 $+18.8\%$、行程 $+18.7\%$ &
\textbf{负优化} & 演练池 30 例对照 \\
N2 & 清除尝试预算 \texttt{max\_attempts} $6\to10$ & 时间 $+5.7\%$，完成率不变 & \textbf{负优化} & 20 例对照 \\
N3 & \texttt{max\_attempts} $6\to16$ & 时间 $+13.3\%$，完成率反而 $-0.003$ & \textbf{负优化} & 同上 \\
N4 & 硬上限 \texttt{hard\_attempt\_cap} $14\to24$ & 时间 $+14.3\%$，完成率不变 & \textbf{负优化} & 同上 \\
N5 & 硬上限 $14\to40$ & 时间 $+20.0\%$，完成率不变 & \textbf{负优化} & 同上 \\
N6 & 末段半径 $90\to130$ m $+$ 硬上限 40 & 时间 $+9.8\%$，完成率 $-0.0005$ & \textbf{负优化} & 同上 \\
N7 & 机会式补测间隔 $650\to450$ m & 完成率 $1.0000\to0.9919$ & \textbf{负优化} & 20 例对照 \\
N8 & 格点间距 $1100\to850$ m（更密） & 完成率 $1.0000\to0.9967$，时间 $+4\%$ & \textbf{负优化} & 20 例对照 \\
N9 & 格点间距 $1100\to1700$ m（更疏） & 完成率 $\to0.9948$ & \textbf{负优化} & 同上 \\
N10 & 全域外环 \texttt{outer\_ring\_gap}$>0$ & 时间 $\times2.4$ & \textbf{负优化} & 30 例对照 \\
N11 & Or-opt 初版（循环内全量重算） & 单局 CPU 数十秒 & \textbf{性能负优化} & 见正文 \\
N12 & C1 顺路清除耦合 & 881.3 s $\to$ 882.2 s（$+0.1\%$） & \textbf{无效} & 24 例对照 \\
N13 & \texttt{progress\_ratio} $0.7\to0.9/0.99$ & 逐位相同 & \textbf{无效} & 20 例对照 \\
\bottomrule
\end{longtable}

\subsection{负优化的共同模式}

复盘之后可以看到三条规律，它们比单个结论更有价值：

\begin{enumerate}[leftmargin=1.8em]
  \item \textbf{"更努力"几乎总是错的}：N2--N6 全是"给机器人更多尝试预算/更大图案"。
        在预算已经被算例检验过的情况下，单纯放宽容忍度只会让它在\textbf{同一个错误估计上}
        多花时间。真正需要的是\textbf{换个地方观测}（垂向截断），而不是\textbf{在原地更努力}。
  \item \textbf{"更密/更疏的采样"两个方向都错}：N8/N9 说明格点间距有一个明确的甜点，
        它的存在是因为"覆盖保证"给了下界（间距须 $\le\sqrt3 R_{\min}\approx1732$ m），
        而"认证成本"给了上界。盲目加密并不提高完成率，因为瓶颈不在覆盖率而在\textbf{末段定位精度}。
  \item \textbf{显式的耦合不一定有增量}：N12 的顺路清除之所以无效，
        是因为任务排序中的优先权偏置已经隐含地实现了同一个效果。
        \textbf{先测量已有机制覆盖了多少，再决定要不要新增机制。}
\end{enumerate}

% ===========================================================================
\section{有效改动清单}

\begin{longtable}{p{3.4cm}p{4.4cm}p{6.2cm}}
\toprule
改动 & 机制 & 实测效果 \\ \midrule
\endhead
增量认证扫描 & 把"未认证位置"的维护改成单调增量，并修正原实现被早退截断的缺陷 &
单局 CPU 41 s $\to$ 2.8 s；同时使规划器看到完整的未认证区域 \\
抑制无法完成认证的搜索站位 & 定向场景下，域内永远无法完成认证的候选点不再吸引站位 &
10 例对照：时间 $-8.5\%$、行程 $-8.7\%$，完成率不变 \\
末段图案中心最后尝试 & 失败后停在估计点而非外圈点 & 位置缺陷修复，代价为零 \\
\textbf{垂向截断测量} & 共线观测无法约束距离，补一次垂向观测 &
$\sigma$ 77 m $\to$ 1 m；案例 1577.9 s $\to$ 739.1 s \\
\textbf{自适应沿射线站位} & 贴边朝外源的受光区是薄月牙，偏移站位落入背光侧 &
两个残留失败算例全部修复 \\
自适应补测间隔 & 未解算频道多时缩短补测间隔 & 24 例对照：时间 $-3.3\%$、行程 $-6.7\%$ \\
坐标下降调参 & 8$\sim$9 个参数的逐维搜索 & 问题三 $-4\%$；问题四见总表 \\
\bottomrule
\end{longtable}

% ===========================================================================
\section{汇总：本轮之后的最终指标}

\begin{center}
\begin{tabular}{lcccc}
\toprule
& 完成率 & 最差完成率 & 平均时间/s & 平均行程/m \\ \midrule
问题三 优化前 & 1.0000 & 1.0000 & 326.3 & 13490 \\
问题三 优化后 & 1.0000 & 1.0000 & 325.8 & 13448 \\
问题四 优化前 & 0.9981 & 0.9231 & 1377.6 & 67501 \\
问题四 优化后 & \textbf{1.0000} & \textbf{1.0000} & \textbf{871.3} & \textbf{41714} \\
\bottomrule
\end{tabular}
\end{center}

在演练池（问题三/四各 60 例）与正式测试上的最终结果：

\begin{center}
\small
\footnotesize
\begin{tabular}{p{4.3cm}p{3.0cm}p{2.8cm}p{4.3cm}}
\toprule
 & 完成率 & 平均时间/s & 说明 \\ \midrule
问题三 演练 60 例 & \textbf{1.0000} & 326.8 & 搜索成本被更多目标分摊 \\
问题三 正式测试 3 次 & 15/15、12/12、10/10 & 257.5 / 392.0 / 389.1 & 均值 346.2 s \\
问题四 演练 60 例（混合） & \textbf{1.0000} & 888.3 & 此前 0.9574 / 661.0 s（错误认证） \\
问题四 演练 30 例（全定向） & \textbf{1.0000} & 983.4 & 此前 0.9164 / 835.2 s \\
问题四 演练 30 例（全全向） & 1.0000 & 860.6 & --- \\
问题四 正式测试 3 次 & \textbf{11/11、12/12、11/11} & 1133 / 945 / 1193 & 三局全部清除 \\
问题四 灵敏度（8 组扰动） & \textbf{全部 1.000} & 711$\sim$1125 & 不依赖单一参数取值 \\
贪心对照（问题四，30 例） & 0.5654 & --- & 无认证判据时大幅下降 \\
\bottomrule
\end{tabular}
\end{center}

（独立算例池：问题三种子 12000--12039，问题四种子 62000--62039，均从未参与调参。）

问题三的变化在噪声内，说明\textbf{问题三已经到结构平台}：8 个参数在 $313\sim348$ s 的窄带内无差异，
瓶颈是"覆盖 $+$ 认证"的固有行程。问题四则由两项定位精度的改动获得决定性改善。

\end{document}
